\section{Bell-polynomial logarithmic calculus}\label{app:bell}

The exponential partial Bell polynomials satisfy
\begin{equation}\label{eq:bell-rec}
\Bell_{n,k}(x_1,\ldots)=
\sum_{j=1}^{n-k+1}\binom{n-1}{j-1}x_j
\Bell_{n-j,k-1},
\quad \Bell_{0,0}=1.
\end{equation}
Fa\`a di Bruno's formula with $f(u)=\log u$ and $u=1+q$ gives
\[
\partial^m\log(1+q)=
\sum_{\ell=1}^m
\frac{(-1)^{\ell-1}(\ell-1)!}{(1+q)^\ell}
\Bell_{m,\ell}(q',\ldots,q^{(m-\ell+1)}).
\]
The Bell polynomials have nonnegative integer coefficients.  If
$|q^{(j)}|\le Q_j$ and $|q|\le Q_0<1$, taking absolute values produces
\eqref{eq:Bm}.  The first three instances are
\begin{align*}
B_1&=\frac{Q_1}{d},\\
B_2&=\frac{Q_2}{d}+\frac{Q_1^2}{d^2},\\
B_3&=\frac{Q_3}{d}+\frac{3Q_1Q_2}{d^2}
+\frac{2Q_1^3}{d^3},\qquad d=1-Q_0.
\end{align*}
Thus the familiar low-order quotient bounds are consequences of one general
recurrence rather than separate calculations.
